\chapter{Metody uczenia przez wzmacnianie}
\label{chap:metody-rl}

W niniejszym rozdziale przedstawiono elementy warstwy uczenia przez
wzmacnianie budowanej nad środowiskiem Ultimate Texas Hold'em.
Pierwszym etapem jest przekształcenie obserwacji udostępnianej przez
silnik gry do stałowymiarowej reprezentacji numerycznej, która może
zostać wykorzystana jako wejście modeli uczących się.

Środowisko udostępnia agentowi obiekt \texttt{UTHObservation}.
Zawiera on wyłącznie informacje dostępne graczowi, w szczególności
jego dwie karty prywatne, odkryte karty wspólne, aktualną fazę
rozdania oraz zbiór legalnych akcji. Reprezentacja numeryczna jest
budowana wyłącznie na podstawie tego obiektu, dzięki czemu proces
uczenia nie uzyskuje dostępu do kart krupiera ani innych informacji
ukrytych w pełnym stanie środowiska.

\section{Reprezentacja stanu}
\label{sec:reprezentacja-stanu}

W eksperymentach przewidziano dwa warianty reprezentacji obserwacji.
Pierwszy wariant stanowi reprezentacja surowa, która koduje informacje
dostępne w \texttt{UTHObservation} bez dodawania ręcznie
zaprojektowanych cech pokerowych. Drugi wariant będzie rozszerzał
reprezentację surową o cechy domenowe wyznaczane na podstawie kart
widocznych dla gracza.

Rozdzielenie reprezentacji od implementacji konkretnych algorytmów
pozwala wykorzystać ten sam sposób kodowania zarówno w agencie
wykorzystującym Deep Q-Network, jak i w agencie opartym na metodzie
Policy Gradient. Dzięki temu wpływ reprezentacji stanu może być
analizowany niezależnie od wybranego algorytmu uczenia.

\subsection{Wspólny interfejs kodowania}
\label{subsec:state-encoder}

Warstwa reprezentacji definiuje protokół \texttt{StateEncoder}.
Każdy encoder udostępnia metodę \texttt{encode()}, która przyjmuje
obiekt \texttt{UTHObservation} i zwraca wektor wartości
zmiennoprzecinkowych. Dodatkowa właściwość \texttt{output\_size}
określa stały wymiar generowanego wektora.

Wspólny kontrakt ma postać

\begin{equation}
    f_{\mathrm{enc}}
    \colon
    O_t
    \rightarrow
    \mathbb{R}^{d},
\end{equation}

gdzie \(O_t\) oznacza obserwację w chwili \(t\), natomiast \(d\)
jest wymiarem wybranej reprezentacji.

Encoder nie zależy od biblioteki wykorzystywanej później do budowy
sieci neuronowych. Wynikiem kodowania jest zwykły wektor numeryczny,
który może zostać następnie przekształcony do formatu wymaganego
przez konkretny algorytm uczenia.

\section{Surowa reprezentacja obserwacji}
\label{sec:raw-state-representation}

Pierwszy wariant reprezentacji został zaimplementowany w klasie
\texttt{RawStateEncoder}. Jego zadaniem jest zachowanie informacji
zawartych bezpośrednio w obserwacji bez wprowadzania jawnych informacji
o sile układu pokerowego, drawach lub strukturze stołu.

Reprezentacja składa się z czterech części:

\begin{enumerate}
    \item zakodowanych dwóch kart prywatnych gracza,
    \item pięciu slotów przeznaczonych na karty wspólne,
    \item kodowania aktualnej fazy rozdania,
    \item maski legalnych akcji.
\end{enumerate}

\subsection{Kodowanie kart}
\label{subsec:kodowanie-kart-raw}

Standardowa talia zawiera 52 unikalne karty. Każda karta jest
reprezentowana przez wektor one-hot o długości 52. Dokładnie jedna
pozycja wektora przyjmuje wartość \(1\), natomiast pozostałe pozycje
mają wartość \(0\).

Karty uporządkowano najpierw według koloru, a następnie według rangi.
Taki porządek jest deterministyczny i pozostaje zgodny ze sposobem
tworzenia standardowej talii wykorzystywanej przez środowisko.

Dla karty \(c\) jej reprezentację można zapisać jako

\begin{equation}
    x_c =
    \left[
        x_0,
        x_1,
        \ldots,
        x_{51}
    \right],
\end{equation}

gdzie

\begin{equation}
    x_i =
    \begin{cases}
        1, & i = \operatorname{index}(c), \\
        0, & i \neq \operatorname{index}(c).
    \end{cases}
\end{equation}

W obserwacji zawsze występują dwie karty prywatne gracza. Dla kart
wspólnych zarezerwowano natomiast pięć pozycji niezależnie od
aktualnej fazy rozdania. Nieodkryte jeszcze karty wspólne są
reprezentowane przez zerowy wektor długości 52.

Takie rozwiązanie zapewnia stały wymiar reprezentacji przed flopem,
po flopie oraz po odkryciu wszystkich kart wspólnych.

\subsection{Kodowanie fazy rozdania}
\label{subsec:kodowanie-fazy-raw}

Uwzględniane są trzy fazy, w których agent może podjąć decyzję:
\texttt{PREFLOP}, \texttt{FLOP} oraz \texttt{RIVER}. Faza jest
reprezentowana przez trzyelementowy wektor one-hot.

Odpowiednie kodowania mają postać

\begin{equation}
    \begin{aligned}
        \mathrm{PREFLOP} &\rightarrow [1,0,0], \\
        \mathrm{FLOP}    &\rightarrow [0,1,0], \\
        \mathrm{RIVER}   &\rightarrow [0,0,1].
    \end{aligned}
\end{equation}

Stan terminalny nie jest kodowany przez \texttt{RawStateEncoder},
ponieważ nie wymaga już podjęcia kolejnej decyzji przez agenta.

\subsection{Maska legalnych akcji}
\label{subsec:maska-akcji-raw}

Ostatnim elementem reprezentacji jest maska legalnych akcji.
Środowisko definiuje sześć możliwych wartości typu \texttt{Action}.
Dla każdej z nich zapisywana jest wartość \(1\), jeżeli akcja jest
legalna w aktualnej fazie, oraz \(0\) w przeciwnym przypadku.

Maska ma więc postać

\begin{equation}
    m_t \in \{0,1\}^{6}.
\end{equation}

Jej obecność pozwala przyszłym agentom jednoznacznie rozpoznać zbiór
akcji dostępnych w danym punkcie decyzyjnym. Mechanizm ten będzie
również wykorzystywany podczas wyboru akcji przez modele uczące się.

\subsection{Wymiar reprezentacji}
\label{subsec:wymiar-raw}

Na kodowanie kart przeznaczono siedem slotów, z czego dwa odpowiadają
kartom gracza, a pięć kartom wspólnym. Każdy slot ma długość 52.
Dodatkowo reprezentacja zawiera trzy wartości opisujące fazę oraz
sześć wartości maski akcji.

Całkowity wymiar reprezentacji wynosi zatem

\begin{equation}
    d_{\mathrm{raw}}
    =
    7 \cdot 52
    + 3
    + 6
    =
    373.
\end{equation}

Wymiar pozostaje stały dla wszystkich stanów decyzyjnych. Jest to
istotne z punktu widzenia późniejszych modeli neuronowych, ponieważ
warstwa wejściowa sieci może posiadać z góry określoną liczbę
elementów.

\subsection{Charakter reprezentacji surowej}
\label{subsec:charakter-raw}

Reprezentacja surowa celowo nie zawiera informacji takich jak
kategoria aktualnego układu, liczba kart do koloru, możliwość
utworzenia strita ani innych cech wynikających z wiedzy pokerowej.
Model otrzymuje jedynie zakodowane dane dostępne bezpośrednio
w obserwacji.

Wariant ten będzie stanowił punkt odniesienia dla reprezentacji
wzbogaconej o cechy domenowe. Porównanie obu wariantów pozwoli ocenić,
czy jawne dostarczenie modelowi informacji pokerowych wpływa na
skuteczność i przebieg procesu uczenia.